\subsection{在数学理论中对组装的应用}

假设集合 S 是某个数学理论 \(\mathfrak{C}\) 的符号集合。我们置 \(n(\square) = 0\) , \(n(\tau) = n(\neg) = 1\) , \(n(\vee) = 2\) , \(n(x) = 0\) 对于每个字母 \(x\) ;

并且最后，对于 \(\mathfrak{C}\) 的每一个特定符号 s， \(n(s)\) 是 s 的权重，当 \(\mathfrak{C}\) 给定时它是固定的。

设 A 是 C 中的一个组装。我们用 \(A^{*}\) 表示通过删除A中的连接所得到的词，并且我们将称A是平衡的，如果 \(A^{*}\) 是平衡的 {[}在 \(L_{0}(S)\) 中{]}。A的一个段是通过如下方式得到的任何组装：在 \(A^{*}\) 的一个段S中，替换在A中连接S中成对符号的那些连接。

\textbf{准则1。} 如果A是T中的一个项或关系，则A是平衡的。

设 \(A_1, A_2, \ldots, A_n\) 是 \(\mathfrak{C}\) 中一个包含 \(A\) 的形成性构造。用归纳法论证，并假设已经证明 \(A_j\) （指标 \(j < i\) ）都是平衡的；现在要证明 \(A_i\) 是平衡的。证明与命题2证明的第一部分完全相同，只有当 \(A_i\) 具有形式 \(\tau_x(B)\) ，其中 \(B = A_j\) 且 \(j < i\) 时例外。在这种情况下，设 \(C\) 为把 \(x\) 在 \(B\) 中的每一处出现替换为 \(\square\) 所得的组装。词 \(A_i^*\) 与 \(\tau C^*\) 相同；现在 \(B^*\) 是平衡的，因而 \(C^*\) 也是平衡的（因为 \(n(\square) = n(x) = 0\) ）。所以 \(A_i^*\) 是平衡的。

我们由此得到了 \(\mathfrak{G}\) 中的一个组装成为一个项或一个关系的必要条件。但这一条件并不充分，我们将会看到。设 \(A\) 是 \(\mathfrak{G}\) 中的一个平衡组装。如果 \(A\) 以一个字母或一个 \(\square\) 开头，那么 \(A\) 必须仅由该符号组成（命题2，推论2）。对于所有其他情形，我们现在将定义 \(A\) 的前导组装。

\begin{enumerate}
\def\labelenumi{(\arabic{enumi})}
\tightlist
\item
  如果 \(A\) 以 \(\lnot\) 开头，或以一个 \(\vee\) 开头，或以一个特定符号开头，则 \(A^*\) 可以恰好以一种方式写成形式 \(fB_1B_2\ldots B_p\) ，其中 \(f\) 是权重 \(p \geqslant 1\) 的符号，并且 \(B_i\) 是平衡的（命题2，推论2）。段 \(A_1, A_2, \ldots, A_p\) （ \(A\) 的段）——它们对应于 \(B_1, B_2, \ldots, B_p\) （ \(A^*\) 的各段） ——被称为 \(A\) 的前导组装。此外，我们将说 \(A\) 完全平衡，如果 \(A\) 与
\end{enumerate}

\[
f A _ {1} A _ {2} \dots A _ {p},
\]

换言之，如果A中没有连接线将f连接到其中一个 \(B_{i}\) ，或连接两个不同的 \(B_{i}\) .

\begin{enumerate}
\def\labelenumi{(\arabic{enumi})}
\setcounter{enumi}{1}
\tightlist
\item
  如果 \(A\) 以 \(\tau\) 开头，那么 \(A^*\) 具有形式 \(\tau B\) ，其中 \(B\) 是平衡的（命题2，推论2）。在这种情况下， \(A\) 的一个前导组装是如下定义的组装 \(A_1\) 中的任何一个：将符号 \(\square\) —— \(B\) 中那些在 \(A\) 里与初始 \(\tau\) 相连接的——替换为一个字母 \(x\) ，该字母不同于 \(B\) 中出现的其他字母，并替换在 \(A\) 中连接 \(B\) 的两个符号的那些连接线。（如果代替 \(x\) ，我们代入一个字母 \(y\) ，它同样不出现在 \(B\) 中，我们得到一个组装，它恰好是 \((y|x)A_1\) 。）此外，我们将说 \(A\) 完全平衡，如果 \(A\) 与 \(\tau_{x}(A_1)\) 相同，换言之，如果没有连接线将初始 \(\tau\) 连接到 \(\pmb{B}\) 的除 \(\square\) 以外的任何符号。
\end{enumerate}

我们现在可以陈述以下准则：

准则2. 设A为T中的平衡组装。

要使A成为一个项，必要且充分的条件是满足以下条件之一：（1）A由单个字母组成；（2）A以 \(\tau\) 开头，是完全平衡的，且其前导组装是关系（由CF8，验证一个前导组装是关系就足够了）；（3）A以实体符号开始，是完全平衡的，且其前导组装是项。要使A是关系，必要且充分的条件是满足以下条件之一：（1）A以 \(\vee\) 或 \(\neg\) 开始，是完全平衡的，且其前导组装是关系；（2）A以关系符号开始，是完全平衡的，且其前导组装是项。

准则CF1至CF4（§1，第4号）表明这些条件是充分的。让我们证明它们是必要的。我们已经在（§1，第3号）中看到，如果A是关系，那么A以∨、或 \(\neg\)、或关系符号开始。在三种情况中的每一种，推理都是类似的。例如，如果A以∨开始，那么A具有∨ BC的形式，其中B和C是关系，因此B和C是A的前导组装；从而A是完全平衡的。如果A是项，那么有两种情况：它由单个字母组成，或者它以实体符号或τ开始。在第二种情况下，我们如上所述进行论证。如果A以τ开始，形成性构造的定义表明A具有τₓ(B)的形式，其中B是关系而X是字母，因此我们可以取B作为A的前导组装，且A是完全平衡的。

当我们想要检验给定的组装A（不是由单个字母组成）在C中是否是关系（相应地，项）时，我们首先验证A是平衡的，并且它以∨、 \(\neg\)、或关系符号开始（相应地，以τ或实体符号开始）。然后我们形成前导组装，并验证（如果适当的话）A是完全平衡的。完成这些之后，我们剩下一个关于较短组装的类似问题。因此，逐步地，我们归结为每个都由单个符号组成的组装，对于这些组装，解答是直接的。

注记。除了某些在公理方面特别弱的数学理论（参见练习7）之外，我们没有这种类型的一般程序来使我们能够检验给定理论G中的给定关系R是否是G中的定理。

